\documentclass[11pt,letterpaper]{article}
\usepackage[top=0.55in, bottom=0.4in, left=0.7in, right=0.7in]{geometry}
\usepackage{lmodern}
\usepackage{hyperref}
\hypersetup{colorlinks=false, urlcolor=black, pdfborder=0 0 0}
\setlength{\parindent}{0pt}
\setlength{\parskip}{0.65em}
\raggedright
\pagestyle{empty}

\begin{document}

\noindent
{\bfseries Ke-Cheng Hsu (Clio)}\\
Ph.D.\ Candidate, Department of Economics, HKUST\\
(852) 6840\,9753 \quad $|$ \quad \href{mailto:khsu@connect.ust.hk}{khsu@connect.ust.hk}\\
\href{https://cliohsu1997.github.io/}{cliohsu1997.github.io} \quad $|$ \quad \href{https://github.com/cliohsu1997}{github.com/cliohsu1997}

\vspace{0.4em}
\noindent
16 November 2026

\vspace{0.3em}
\noindent
Hiring Team\\
Beevest Financial Trading Limited\\
Re: Data Strategy Analyst

\vspace{0.15em}
Dear Hiring Team,

I am writing to apply for the Data Strategy Analyst position. I am a Ph.D.\ candidate in Economics at HKUST (expected June 2027; I can start as early as February 2027), a Hong Kong permanent resident with the right to work in Hong Kong, a native Mandarin speaker, and professionally fluent in English. The role calls for turning multi-source data into strategy-relevant insights, monitoring portfolio positions and market conditions, supporting valuation and status reporting, and helping with fund-operations work. That matches how I already work: learn a new market or dataset quickly, build a careful analysis pipeline, check the numbers, and explain the result so a non-specialist can act on it.

Academic excellence and fast learning are the core of my profile. I completed a B.Sc.\ in Mathematics and Economics with a minor in Actuarial Mathematics at HKUST (CGA 3.9/4.3; Dean's List; Academic Achievement Medal, top 1\%), an M.Phil.\ in Economics (CGA 4.1/4.3), and I am finishing a Ph.D.\ in empirical industrial organization and digital economics. Across these degrees I repeatedly picked up new tools---structural estimation, GMM, large-scale Python / R / C++ pipelines on AWS and HPC---and used them on real proprietary data without a long onboarding period. I use AI tools daily to accelerate coding and research, then verify outputs with tests so speed never replaces accuracy.

My projects map directly to analysis, monitoring, and reporting. In my sole-authored job market paper on Taobao apparel return policy, I integrate orders, prices, insurance, and returns; identify patterns in purchase and return behavior; and write clear counterfactual comparisons: a menu with a moderate restocking fee raises social welfare by about 2 million RMB (roughly a 5\% gain) and seller profit by about 1.2 million RMB relative to status-quo policies. With Professors Kohei Kawaguchi and Kei Kawai, I monitor pricing, inventory, and return flows in bookstore point-of-sale panels and deliver weekly status-style figures and short written updates on a fixed deadline. Earlier work on an interbank crisis network gave me a concrete frame for positions, exposures, and cascade risk---useful when thinking about portfolio and market conditions. I do not claim years of sell-side derivatives valuation experience; I do bring strong quantitative modelling, careful documentation, and the ability to learn product and valuation conventions quickly on the job.

Clear communication is the other half of the work. As a lecturer and teaching assistant at HKUST (4.53/5 and 4.72/5), and as Teaching Assistant Coordinator in the Business School, I translate technical material into plain language, keep reports consistent before release, and coordinate deliverables across people. I am comfortable preparing periodic status reports, presenting insights in Mandarin and English, and supporting administrative fund-operations tasks with the same attention to detail.

My present monthly compensation is HKD\,20,000 (postgraduate studentship). My expected monthly basic salary is HKD\,25,000--30,000. I attach my resume. Thank you for your consideration.

\vspace{0.4em}
Sincerely,\\[0.35em]
Ke-Cheng Hsu (Clio)

\end{document}
